\section{人の位置情報を活用する研究}
人の位置情報は，行動分析や環境最適化など利用目的に応じて多様な分野で応用が検討されている．
特に屋内外における人の移動や滞在の履歴の取得により，人の行動特性や空間利用の実態を定量的に把握できる点が注目されている．
例えば，Google Maps\cite{GoogleMap}に代表されるナビゲーションシステムでは，位置情報を用いた現在地の可視化と経路案内が広く普及している．
また，位置情報とゲーミフィケーションを組み合わせた Pokémon GO\cite{PokemonGo}のようなゲームでは，実空間での移動をゲーム体験と結び付けユーザの行動変容を促す仕組みが実用化されている．
これらは人の移動や滞在といった行動を位置情報から捉え，その特徴を分析・活用している点に共通点がある．
こういった，人の行動特性を明らかにする行動分析に関する研究が数多く行われている．


行動分析の分野では，位置履歴を用いてエリアごとの滞在時間や訪問頻度を分析し，人の行動傾向や利用パターンを明らかにする研究が数多く行われている．
大規模商業施設において来訪者の歩行軌跡を分析し，空間の実際の利用状況と設計意図との差異を示した研究\cite{layout}が報告されている．
さらに，長期間の位置情報の蓄積から個人や集団の行動パターンを抽出し，行動の類型化や特徴量の分析を行う試みもある．
例えば，人物動線データ群から行動パターンを分類し，通常行動から逸脱した行動を検出する研究\cite{pattern}が報告されている．
また，大地震時を想定した広域避難シミュレーションにおいて，避難軌跡をクラスタリングして行動傾向を把握し，将来の避難行動予測へ応用する研究\cite{simulation}も提案されている．
このように，人の位置情報を時系列に蓄積・分析すると，日常行動から非常時の避難行動に至るまで様々な状況における行動傾向や空間利用の特徴を明らかにできると示されており，位置情報に基づく行動分析は幅広い分野で有効である．


また，人の位置情報は利用目的や環境に応じて多様な形で活用可能であり，本研究で前提となる部屋単位の滞在管理もこうした屋内位置推定の技術における応用の一例として位置づけられる．
滞在管理は，人物がいつどの空間に存在しているかの把握を目的としており，詳細な位置や移動経路を高精度に推定する必要がない点に特徴がある．
特に，オフィスや研究室，大学施設などの屋内環境では，在室状況の把握が円滑な業務遂行やコミュニケーションの促進に寄与すると考えられており，前述した来訪促進を目的とする研究もこの一例である．
この情報は在室人数の把握による空間利用状況の可視化，利用頻度や滞在時間の分析に基づく運用改善，在室情報の共有による対面コミュニケーションの支援などを通じて活用される．


さらに，こうして得られた人の位置や滞在状況の情報は，環境最適化を目的とした研究にも応用されている．
位置情報を用いて空調，照明，案内表示などの制御を行うシステムが提案されており，ユーザの存在や分布に応じて環境を動的に調整する試みが行われている．
例えば，滞在状況に応じて照明を制御する知的照明システム\cite{light}や，人の分布に基づいて空調を制御する手法\cite{airControl}などが検討されている．
これらは快適性の向上やエネルギー消費の削減を目的としている．
このように，人の位置情報は分析，最適化，可視化，娯楽など，多様な目的で活用可能な情報であり，屋内外を問わず幅広い応用が検討されている．